\documentclass[11pt,a4paper]{article}
\usepackage[left=24mm,right=24mm,top=22mm,bottom=20mm,includefoot]{geometry}
\usepackage{amsmath,amssymb,amsthm}
\usepackage{iftex}
\ifPDFTeX
  \usepackage{CJKutf8}
  % CJKutf8 supplies its own Unicode map for each Chinese subfont.
  \makeatletter
  \AddToHook{cmd/CJK@addcmap/before}{\pdfnobuiltintounicode\font@name}
  \makeatother
\else
  \usepackage[no-math]{fontspec}
  \IfFontExistsTF{Noto Serif CJK SC}{
    \setmainfont{Noto Serif CJK SC}
    \setsansfont{Noto Sans CJK SC}
  }{
    \setmainfont{FandolSong-Regular.otf}[BoldFont=FandolSong-Bold.otf]
    \setsansfont{FandolHei-Regular.otf}[BoldFont=FandolHei-Bold.otf]
  }
  \ifXeTeX
    \XeTeXlinebreaklocale "zh"
    \XeTeXlinebreakskip=0pt plus 0.5pt
  \fi
\fi
\pagestyle{plain}
\raggedbottom
\linespread{1.13}
\setlength{\parindent}{2em}
\setlength{\parskip}{3pt plus 1pt}
\setlength{\jot}{3pt}
\DeclareMathSizes{10.95}{10.95}{9}{7}
% Use a single stroke for Chinese headings when pdfLaTeX has no bold CJK font.
\newcommand{\headingweight}{\ifPDFTeX\mdseries\else\bfseries\fi}
\makeatletter
\renewcommand{\section}{\@startsection{section}{1}{0pt}
  {2.0ex plus .3ex minus .2ex}{.9ex}{\normalfont\large\headingweight}}
\makeatother

\newcommand{\indep}{\mathrel{\perp\!\!\!\perp}}
\newtheoremstyle{report}{7pt}{7pt}{\normalfont}{}{\headingweight}{。}{0.4em}{}
\theoremstyle{report}
\newtheorem{assumption}{假设}
\newtheorem{theorem}{定理}
\begin{document}
\ifPDFTeX\begin{CJK*}{UTF8}{gbsn}\fi
\begin{center}
{\LARGE\headingweight 两个暴露下有序结局 Shapley 责任的识别}\\[5pt]
{\large 技术报告}
\end{center}

设 $X_1,X_2\in\{0,1\}$，$Y\in\{0,1,2\}$，数值越大表示结局越严重。
四个潜在结果组成类型 $T=(Y_{00},Y_{01},Y_{10},Y_{11})$。
按每升一级计 $1$ 分，定义病例事件
\[
\mathcal E_1=\{X_1=X_2=1,Y=1\},\qquad
\mathcal E_2=\{X_1=X_2=1,Y=2\}.
\]
对满足 $t_{11}=2$ 的类型 $t$，将各潜在结果取为该类型的相应分量，定义
\begin{equation}
\begin{aligned}
R(1\mid t,\mathcal E_2)&=\tfrac12\{(Y_{10}-Y_{00})+(Y_{11}-Y_{01})\},\\
R(2\mid t,\mathcal E_2)&=\tfrac12\{(Y_{01}-Y_{00})+(Y_{11}-Y_{10})\},\\
R(0\mid t,\mathcal E_2)&=Y_{00}.
\end{aligned}
\label{eq:responsibility}
\end{equation}
其中 $0$ 表示无暴露背景。对满足 $t_{11}=1$ 的类型，$R(j\mid t,\mathcal E_1)$
由同一规则定义。对 $j=0,1,2$，病例平均责任为
\begin{equation}
\begin{aligned}
\mathcal R_j(\mathcal E_1)
 &=\sum_{t:\,t_{11}=1}R(j\mid t,\mathcal E_1)\,
     \mathrm{pr}(T=t\mid\mathcal E_1),\\
\mathcal R_j(\mathcal E_2)
 &=\sum_{t:\,t_{11}=2}R(j\mid t,\mathcal E_2)\,
     \mathrm{pr}(T=t\mid\mathcal E_2).
\end{aligned}
\label{eq:target}
\end{equation}

\section{假设}
\begin{assumption}[一致性、随机分配与正值性]\label{ass:observed}
当 $(X_1,X_2)$ 分别为 $(0,0),(0,1),(1,0),(1,1)$ 时，
观测结局分别等于 $Y_{00},Y_{01},Y_{10},Y_{11}$。此外，
\[
(X_1,X_2)\indep T,\qquad
\mathrm{pr}(X_1=a,X_2=b)>0,\quad a,b=0,1.
\]
\end{assumption}

\begin{assumption}[单步单调增量]\label{ass:unit}
下列不等式几乎处处成立：
\begin{equation}
\begin{aligned}
0\le Y_{10}-Y_{00}\le1,\qquad&0\le Y_{11}-Y_{10}\le1,\\
0\le Y_{01}-Y_{00}\le1,\qquad&0\le Y_{11}-Y_{01}\le1.
\end{aligned}
\label{eq:unit}
\end{equation}
\end{assumption}

\begin{assumption}[Markov 条件]\label{ass:markov}
同一个潜在结果联合分布满足
\begin{equation}
Y_{11}\indep Y_{00}\mid Y_{10},
\qquad
Y_{11}\indep Y_{00}\mid Y_{01}.
\label{eq:markov}
\end{equation}
\end{assumption}

\clearpage
\section{定理}
记
\[
\begin{aligned}
\xi_{00}&=\mathrm{pr}(Y_{00}=2),&\qquad
\xi_{01}&=\mathrm{pr}(Y_{01}=2),\\
\xi_{10}&=\mathrm{pr}(Y_{10}=2),&\qquad
\xi_{11}&=\mathrm{pr}(Y_{11}=2),
\end{aligned}
\]
以及 $\theta=\mathrm{pr}(Y_{00}=0,Y_{11}=2)$。

\begin{theorem}[病例平均责任的点识别]
在假设~\ref{ass:observed}--\ref{ass:markov} 下，
$\mathcal R_j(\mathcal E_1)$ 和 $\mathcal R_j(\mathcal E_2)$
在相应病例事件发生概率为正时均可由观测分布点识别。
其中
\begin{equation}
\theta=
\begin{cases}
\dfrac{\{\mathrm{pr}(Y_{00}=0)-\mathrm{pr}(Y_{10}=0)\}
       (\xi_{11}-\xi_{10})}
      {\mathrm{pr}(Y_{10}=1)},
   & \mathrm{pr}(Y_{10}=1)>0,\\[10pt]
0, & \mathrm{pr}(Y_{10}=1)=0.
\end{cases}
\label{eq:theta}
\end{equation}
当 $\xi_{11}>0$ 时，
\begin{equation}
\begin{aligned}
\mathcal R_1(\mathcal E_2)
 &=\frac{\xi_{11}-\xi_{01}+\xi_{10}-\xi_{00}+\theta}{2\xi_{11}},\\[4pt]
\mathcal R_2(\mathcal E_2)
 &=\frac{\xi_{11}-\xi_{10}+\xi_{01}-\xi_{00}+\theta}{2\xi_{11}},\\[4pt]
\mathcal R_0(\mathcal E_2)
 &=\frac{\xi_{11}+\xi_{00}-\theta}{\xi_{11}}.
\end{aligned}
\label{eq:grade2}
\end{equation}
当 $\mathrm{pr}(Y_{11}=1)>0$ 时，
\begin{equation}
\begin{aligned}
\mathcal R_1(\mathcal E_1)
 &=\frac{\mathrm E\{(Y_{10}-Y_{00})+(Y_{11}-Y_{01})\}
       -\xi_{11}+\xi_{01}-\xi_{10}+\xi_{00}-\theta}
      {2\mathrm{pr}(Y_{11}=1)},\\[4pt]
\mathcal R_2(\mathcal E_1)
 &=\frac{\mathrm E\{(Y_{01}-Y_{00})+(Y_{11}-Y_{10})\}
       -\xi_{11}+\xi_{10}-\xi_{01}+\xi_{00}-\theta}
      {2\mathrm{pr}(Y_{11}=1)},\\[4pt]
\mathcal R_0(\mathcal E_1)
 &=\frac{\mathrm E(Y_{00})-\xi_{11}-\xi_{00}+\theta}
      {\mathrm{pr}(Y_{11}=1)}.
\end{aligned}
\label{eq:grade1}
\end{equation}
上述期望均为无条件期望；其被积函数是四个潜在结果的线性组合。
对发生概率为正的相应事件，
\[
\sum_{j=0}^{2}\mathcal R_j(\mathcal E_1)=1,\qquad
\sum_{j=0}^{2}\mathcal R_j(\mathcal E_2)=2.
\]
\end{theorem}

\clearpage
\section{证明}
由假设~\ref{ass:observed}，四个潜在结果的边际分布分别由四个暴露组的结局分布识别。
一致性与联合随机分配还给出
\[
\begin{aligned}
\mathrm{pr}(T=t\mid\mathcal E_2)
 &=\frac{\mathrm{pr}(T=t,X_1=X_2=1,Y_{11}=2)}
         {\mathrm{pr}(X_1=X_2=1,Y_{11}=2)}\\
 &=\frac{\mathrm{pr}(T=t,Y_{11}=2)}{\mathrm{pr}(Y_{11}=2)}
 =\mathrm{pr}(T=t\mid Y_{11}=2).
\end{aligned}
\]
$\mathcal E_1$ 同理。因此，式~\eqref{eq:target} 中的类型加权平均，
分别等于式~\eqref{eq:responsibility} 右侧给定 $Y_{11}=1$ 或 $Y_{11}=2$ 的条件期望。

假设~\ref{ass:unit} 蕴含
\[
Y_{11}-1\le Y_{10}\le Y_{00}+1,\qquad
Y_{11}-1\le Y_{01}\le Y_{00}+1.
\]
当 $Y_{00}=0,Y_{11}=2$ 时，上述上下界均为 $1$，
故 $Y_{10}=Y_{01}=1$，从而
\begin{equation}
\theta=\mathrm{pr}(Y_{00}=0,Y_{10}=1,Y_{11}=2).
\label{eq:theta-event}
\end{equation}

由 $0\le Y_{10}-Y_{00}\le1$，在 $Y_{00}=0$ 时，$Y_{10}$ 只可能取 $0,1$；
而 $Y_{10}=0$ 必有 $Y_{00}=0$。因此
\[
\begin{aligned}
\mathrm{pr}(Y_{00}=0)
 &=\mathrm{pr}(Y_{00}=0,Y_{10}=0)
   +\mathrm{pr}(Y_{00}=0,Y_{10}=1)\\
 &=\mathrm{pr}(Y_{10}=0)+\mathrm{pr}(Y_{00}=0,Y_{10}=1).
\end{aligned}
\]
同理，$0\le Y_{11}-Y_{10}\le1$ 给出
\[
\xi_{11}
=\mathrm{pr}(Y_{10}=1,Y_{11}=2)
 +\mathrm{pr}(Y_{10}=2,Y_{11}=2)
=\mathrm{pr}(Y_{10}=1,Y_{11}=2)+\xi_{10}.
\]
故
\begin{equation}
\begin{aligned}
\mathrm{pr}(Y_{00}=0,Y_{10}=1)
 &=\mathrm{pr}(Y_{00}=0)-\mathrm{pr}(Y_{10}=0),\\
\mathrm{pr}(Y_{10}=1,Y_{11}=2)&=\xi_{11}-\xi_{10},\\
\mathrm{pr}(Y_{01}=1,Y_{11}=2)&=\xi_{11}-\xi_{01}.
\end{aligned}
\label{eq:adjacent}
\end{equation}
第三式由 $0\le Y_{11}-Y_{01}\le1$ 以相同方式得到。

当 $\mathrm{pr}(Y_{10}=1)>0$ 时，式~\eqref{eq:theta-event} 与假设~\ref{ass:markov} 给出
\begin{equation}
\begin{aligned}
\theta
 &=\mathrm{pr}(Y_{10}=1)\,
   \mathrm{pr}(Y_{00}=0,Y_{11}=2\mid Y_{10}=1)\\
 &\overset{\text{假设~\ref{ass:markov}}}{=}
   \mathrm{pr}(Y_{10}=1)\,
   \mathrm{pr}(Y_{00}=0\mid Y_{10}=1)\,
   \mathrm{pr}(Y_{11}=2\mid Y_{10}=1)\\
 &=\frac{\mathrm{pr}(Y_{00}=0,Y_{10}=1)\,
           \mathrm{pr}(Y_{10}=1,Y_{11}=2)}
          {\mathrm{pr}(Y_{10}=1)}\\
 &=\frac{\{\mathrm{pr}(Y_{00}=0)-\mathrm{pr}(Y_{10}=0)\}
           (\xi_{11}-\xi_{10})}
          {\mathrm{pr}(Y_{10}=1)}.
\end{aligned}
\label{eq:theta-factor}
\end{equation}
若 $\mathrm{pr}(Y_{10}=1)=0$，由式~\eqref{eq:theta-event} 有 $\theta=0$。
这证明了式~\eqref{eq:theta}。

\clearpage
由单调性，$\mathrm{pr}(Y_{10}=2,Y_{11}=2)=\xi_{10}$，
$\mathrm{pr}(Y_{01}=2,Y_{11}=2)=\xi_{01}$。结合式~\eqref{eq:adjacent}，
\begin{equation}
\begin{aligned}
\mathrm E\{Y_{10}\mathbf1(Y_{11}=2)\}
 &=\mathrm{pr}(Y_{10}=1,Y_{11}=2)
     +2\mathrm{pr}(Y_{10}=2,Y_{11}=2)\\
 &=(\xi_{11}-\xi_{10})+2\xi_{10}
  =\xi_{11}+\xi_{10},\\[3pt]
\mathrm E\{Y_{01}\mathbf1(Y_{11}=2)\}
 &=\mathrm{pr}(Y_{01}=1,Y_{11}=2)
     +2\mathrm{pr}(Y_{01}=2,Y_{11}=2)\\
 &=(\xi_{11}-\xi_{01})+2\xi_{01}
  =\xi_{11}+\xi_{01}.
\end{aligned}
\label{eq:exposure-moments}
\end{equation}

背景的三个可能取值将事件 $Y_{11}=2$ 分解为互不相交的三部分：
\[
\begin{aligned}
\xi_{11}
 &=\mathrm{pr}(Y_{00}=0,Y_{11}=2)
   +\mathrm{pr}(Y_{00}=1,Y_{11}=2)
   +\mathrm{pr}(Y_{00}=2,Y_{11}=2)\\
 &=\theta+\mathrm{pr}(Y_{00}=1,Y_{11}=2)+\xi_{00}.
\end{aligned}
\]
因此
\begin{equation}
\begin{aligned}
\mathrm E\{Y_{00}\mathbf1(Y_{11}=2)\}
 &=\mathrm{pr}(Y_{00}=1,Y_{11}=2)
     +2\mathrm{pr}(Y_{00}=2,Y_{11}=2)\\
 &=(\xi_{11}-\theta-\xi_{00})+2\xi_{00}
  =\xi_{11}+\xi_{00}-\theta.
\end{aligned}
\label{eq:background-moment}
\end{equation}
此外，$\mathrm E\{Y_{11}\mathbf1(Y_{11}=2)\}=2\xi_{11}$。

对式~\eqref{eq:responsibility} 取带指标的无条件期望，
由式~\eqref{eq:exposure-moments}--\eqref{eq:background-moment} 得
\begin{equation}
\begin{aligned}
&\mathrm E\!\left[\{(Y_{10}-Y_{00})+(Y_{11}-Y_{01})\}
                      \mathbf1(Y_{11}=2)\right]\\
&\quad=(\xi_{11}+\xi_{10})-(\xi_{11}+\xi_{00}-\theta)
          +2\xi_{11}-(\xi_{11}+\xi_{01})\\
&\quad=\xi_{11}-\xi_{01}+\xi_{10}-\xi_{00}+\theta,\\[4pt]
&\mathrm E\!\left[\{(Y_{01}-Y_{00})+(Y_{11}-Y_{10})\}
                      \mathbf1(Y_{11}=2)\right]\\
&\quad=(\xi_{11}+\xi_{01})-(\xi_{11}+\xi_{00}-\theta)
          +2\xi_{11}-(\xi_{11}+\xi_{10})\\
&\quad=\xi_{11}-\xi_{10}+\xi_{01}-\xi_{00}+\theta.
\end{aligned}
\label{eq:grade2-numerators}
\end{equation}
当 $\xi_{11}>0$ 时，
\[
\begin{aligned}
\mathcal R_1(\mathcal E_2)
 &=\frac{\mathrm E[\{(Y_{10}-Y_{00})+(Y_{11}-Y_{01})\}
                          \mathbf1(Y_{11}=2)]}{2\xi_{11}},\\
\mathcal R_2(\mathcal E_2)
 &=\frac{\mathrm E[\{(Y_{01}-Y_{00})+(Y_{11}-Y_{10})\}
                          \mathbf1(Y_{11}=2)]}{2\xi_{11}},\\
\mathcal R_0(\mathcal E_2)
 &=\frac{\mathrm E\{Y_{00}\mathbf1(Y_{11}=2)\}}{\xi_{11}}.
\end{aligned}
\]
将式~\eqref{eq:grade2-numerators} 与式~\eqref{eq:background-moment} 代入，得到式~\eqref{eq:grade2}。

\clearpage
假设~\ref{ass:unit} 蕴含 $Y_{11}=0$ 时 $Y_{00}=Y_{01}=Y_{10}=0$。
因而在此事件上，
\[
(Y_{10}-Y_{00})+(Y_{11}-Y_{01})=0,\qquad
(Y_{01}-Y_{00})+(Y_{11}-Y_{10})=0,\qquad Y_{00}=0.
\]
又因 $Y_{11}$ 只取 $0,1,2$，有以下逐个体恒等式：
\begin{equation}
\begin{aligned}
&\{(Y_{10}-Y_{00})+(Y_{11}-Y_{01})\}\mathbf1(Y_{11}=1)\\
&\quad=\{(Y_{10}-Y_{00})+(Y_{11}-Y_{01})\}
       -\{(Y_{10}-Y_{00})+(Y_{11}-Y_{01})\}\mathbf1(Y_{11}=2),\\[4pt]
&\{(Y_{01}-Y_{00})+(Y_{11}-Y_{10})\}\mathbf1(Y_{11}=1)\\
&\quad=\{(Y_{01}-Y_{00})+(Y_{11}-Y_{10})\}
       -\{(Y_{01}-Y_{00})+(Y_{11}-Y_{10})\}\mathbf1(Y_{11}=2),\\[4pt]
&Y_{00}\mathbf1(Y_{11}=1)=Y_{00}-Y_{00}\mathbf1(Y_{11}=2).
\end{aligned}
\label{eq:grade1-decomposition}
\end{equation}
取期望并代入式~\eqref{eq:grade2-numerators} 与式~\eqref{eq:background-moment}，得
\begin{equation}
\begin{aligned}
&\mathrm E\!\left[\{(Y_{10}-Y_{00})+(Y_{11}-Y_{01})\}
                       \mathbf1(Y_{11}=1)\right]\\
&\quad=\mathrm E\{(Y_{10}-Y_{00})+(Y_{11}-Y_{01})\}
       -\xi_{11}+\xi_{01}-\xi_{10}+\xi_{00}-\theta,\\[4pt]
&\mathrm E\!\left[\{(Y_{01}-Y_{00})+(Y_{11}-Y_{10})\}
                       \mathbf1(Y_{11}=1)\right]\\
&\quad=\mathrm E\{(Y_{01}-Y_{00})+(Y_{11}-Y_{10})\}
       -\xi_{11}+\xi_{10}-\xi_{01}+\xi_{00}-\theta,\\[4pt]
&\mathrm E\{Y_{00}\mathbf1(Y_{11}=1)\}
   =\mathrm E(Y_{00})-\xi_{11}-\xi_{00}+\theta.
\end{aligned}
\label{eq:grade1-numerators}
\end{equation}
当 $\mathrm{pr}(Y_{11}=1)>0$ 时，式~\eqref{eq:grade1-numerators} 中前两个期望
各除以 $2\mathrm{pr}(Y_{11}=1)$，第三个期望除以 $\mathrm{pr}(Y_{11}=1)$，
分别得到 $\mathcal R_1(\mathcal E_1),\mathcal R_2(\mathcal E_1),\mathcal R_0(\mathcal E_1)$，
即式~\eqref{eq:grade1}。上述计算未除以 $\xi_{11}$，故也适用于 $\xi_{11}=0$。

式~\eqref{eq:theta} 只依赖四个潜在结果的边际概率。
式~\eqref{eq:grade1} 中的无条件期望由线性性写为边际均值之和，例如
\[
\mathrm E\{(Y_{10}-Y_{00})+(Y_{11}-Y_{01})\}
=\mathrm E(Y_{10})-\mathrm E(Y_{00})+\mathrm E(Y_{11})-\mathrm E(Y_{01}).
\]
因此，式~\eqref{eq:theta}--\eqref{eq:grade1} 均为观测分布的函数。

最后，逐个体有
\[
\tfrac12\{(Y_{10}-Y_{00})+(Y_{11}-Y_{01})\}
+\tfrac12\{(Y_{01}-Y_{00})+(Y_{11}-Y_{10})\}
+Y_{00}=Y_{11}.
\]
分别取给定 $Y_{11}=1$ 和 $Y_{11}=2$ 的条件期望，得到定理中的两条责任加总等式。
\hfill$\square$

\ifPDFTeX\end{CJK*}\fi
\end{document}
